\documentclass[11pt]{article}
\usepackage[margin=1in]{geometry}
\usepackage{hyperref}
\usepackage{enumitem}
\usepackage{booktabs}
\usepackage{longtable}
\usepackage{xcolor}
\title{ProbOS --- Overall Plan}
\author{Nisong Monyimba}
\date{\today}

\begin{document}
\maketitle

\section*{Purpose}
This document is the single source of truth for the full ProbOS roadmap.
Earlier planning used loose ``Year 1 / Year 2 / Year 3'' labels that did
not line up with actual calendar years (Year 1 covered 4 months, Year 2
covered 8 months). This document replaces those labels with a
consistent Month 1 through Month 24 numbering. No month number here is
retroactively changed once work on it begins --- if a month's plan turns
out to be wrong once we reach it, we correct it in that month's own
\texttt{main.tex} (as already done for Month 1 and Month 2), not here.

\section*{Total planned horizon}
\textbf{24 months (2 years).} Months 1--12 are concretely planned at
weekly granularity. Months 13--24 are directional and will be replanned
in detail as each one approaches, the same way Month 2 was just
replanned at the end of Month 1.

\section*{Status legend}
\begin{itemize}[leftmargin=*]
  \item \textbf{DONE} --- complete, documented in its own monthly plan
    PDF under \texttt{docs/monthly\_plans/}
  \item \textbf{PLANNED} --- detailed week-by-week plan exists
  \item \textbf{DIRECTIONAL} --- theme and goal known, not yet
    week-by-week planned
\end{itemize}

\section*{Months 1--4 --- Kernel Foundation}

\begin{longtable}{p{1.3cm}p{4.5cm}p{7.5cm}p{1.6cm}}
\toprule
\textbf{Month} & \textbf{Theme} & \textbf{Key deliverables} & \textbf{Status} \\
\midrule
\endhead
1 & Distribution + Model ABCs, Monte Carlo kernel, PDSL v0.1 &
  \texttt{Distribution}/\texttt{Model} ABCs, \texttt{BatteryModel2Cell},
  \texttt{MonteCarloEngine}, \texttt{SobolSensitivity},
  \texttt{ProvenanceTracker}, C++/OpenMP kernel (7x), PDSL compiler v0.1,
  3 cross-discipline examples, 264 Python + 13 C++ tests &
  DONE \\
2 & Sequential inference + C++ into Python package &
  \texttt{ParticleFilter} (validated vs Kalman ground truth), pybind11
  bindings, FastAPI service layer, cross-discipline filtering examples &
  PLANNED \\
3 & GPU kernel path, pharma domain, automotive positioning, PDSL v0.2,
  comprehensive C++/kernel sweep &
  GPU kernel path (honestly benchmarked), \texttt{PharmaStabilityModel}
  (validated against Gonzalez-Gonzalez et al. 2023), automotive/EV
  positioning materials, PDSL v0.2 (control flow), C++/Python
  parameter-spread discrepancy fixed at the root, 408 Python + 5 C++
  test executables &
  DONE \\
3.5 & Injection molding process qualification (Week 14, research-gate
  first) &
  \texttt{InjectionMoldingProcessModel}, reusing the existing engine
  for a third domain; real validation contingent on a genuine Day 1
  research gate; fourth live demo tab &
  PLANNED \\
4 & Dashboard architecture and core build &
  Framework decision (Streamlit vs. React), core dashboard for the
  validated domains, P05/P50/P95 visualization, Sobol chart, error
  handling and a real cold-start user-flow test &
  PLANNED \\
\bottomrule
\end{longtable}

\section*{Months 5--12 --- Directional (formerly ``Year 2'')}

\begin{longtable}{p{1.3cm}p{11.9cm}}
\toprule
\textbf{Month} & \textbf{Directional theme} \\
\midrule
\endhead
5--8 & Consolidation of Months 1--4 capabilities into a stable,
  documented, installable package. Expanded example library across more
  domains. Performance work informed by real profiling from earlier
  months, not assumed in advance. \\
9--12 & Broader platform concerns: packaging for external users,
  documentation, and any capabilities that Months 1--8 reveal are
  actually needed (rather than pre-specified now). \\
\bottomrule
\end{longtable}

\textbf{Honesty note:} the specific technical content of Months 5--12
was previously drafted in detail before Month 1 was even finished. That
level of pre-specification was premature --- Month 1 alone surfaced
priorities (e.g. the particle filter becoming Month 2's centrepiece)
that were not obvious in advance. Months 5--12 are intentionally left
as themes here and will be planned in full, in their own
\texttt{docs/monthly\_plans/monthN/main.tex}, one month at a time as we
approach them.

\section*{Months 13--24 --- Directional (formerly ``Year 3'')}

\begin{longtable}{p{1.3cm}p{11.9cm}}
\toprule
\textbf{Month} & \textbf{Directional theme} \\
\midrule
\endhead
13--18 & Scale and robustness of whatever platform shape has emerged
  from Months 1--12. \\
19--24 & Maturity, external usage, and closing out the 24-month
  horizon with a clear-eyed assessment of what should come next (a new
  24-month plan, or a narrower maintenance phase). \\
\bottomrule
\end{longtable}

\section*{Process going forward}
\begin{enumerate}[leftmargin=*]
  \item This document (\texttt{docs/monthly\_plans/overall/main.tex}) is
    updated whenever the total horizon or month numbering changes ---
    not on every month's completion.
  \item Each individual month gets its own
    \texttt{docs/monthly\_plans/monthN/main.tex} plus compiled PDF,
    written and compiled \emph{before} that month's work begins.
  \item Each month closes with a retrospective appended to (or replacing
    the plan content of) that same month's PDF, as already done for
    Month 1.
  \item \texttt{check\_ci.sh} is run after every push; mypy strict and
    ruff stay clean; no benchmark claim is made without a genuine
    measurement backing it.
\end{enumerate}

\end{document}
